\documentclass[11pt]{article}

\usepackage[margin=1in]{geometry}
\usepackage{amsmath,amssymb,amsthm}
\usepackage{enumitem}
\usepackage[hidelinks]{hyperref}

\newtheorem{theorem}{Theorem}
\newtheorem{proposition}[theorem]{Proposition}
\newtheorem{corollary}[theorem]{Corollary}
\newtheorem{lemma}[theorem]{Lemma}
\theoremstyle{remark}
\newtheorem{remark}[theorem]{Remark}
\newtheorem{question}[theorem]{Question}

\newcommand{\Z}{\mathbb{Z}}
\newcommand{\F}{\mathbb{F}}
\newcommand{\Q}{\mathbb{Q}}
\newcommand{\Cl}{\operatorname{Cl}}
\newcommand{\Sym}{\operatorname{Sym}}
\newcommand{\aw}{\operatorname{aw}}
\newcommand{\mean}{\operatorname{mean}}
\newcommand{\rad}{\operatorname{rad}}
\newcommand{\taglabel}[1]{\textnormal{\textbf{[#1]}}}
\newcommand{\PROVED}{\taglabel{PROVED}}
\newcommand{\TESTED}{\taglabel{TESTED}}
\newcommand{\SOURCE}{\taglabel{SOURCE-PINNED}}
\newcommand{\OPEN}{\taglabel{OPEN}}
\newcommand{\RESOLVED}{\taglabel{RESOLVED}}

\title{Quadratic deformations of the game exterior algebra}
\author{Codex}
\date{2026-06-17}

\begin{document}
\maketitle

\noindent\RESOLVED{} \textbf{Closure addendum (2026-08-09).}
The local coefficient analysis preserved in this note has now been superseded
at the global boundary by \path{writeups/game_exterior_divisibility.tex}.
Moews's structure theorem makes multiplication by every power of two
surjective on the full short-game group.  If a torsion game $t$ has order
$n=2^k$, choose $ny=t$ and $nz=x$; then $n^2y=0$, so every additive grade-one
realization satisfies
\[
  e_t^2=(ne_y)^2=0,
  \qquad e_te_x+e_xe_t=0.
\]
Coefficient faithfulness therefore forces $Q(t)=B(t,x)=0$ in every
coefficient characteristic.  Thus no nonzero torsion table extends as an
ambient-coherent game-native Clifford deformation.  The root-incomplete local
tables studied below remain valid as local algebra, but are not residual cases
of the resolved global question.

\begin{abstract}
This note preserves an earlier attack on the then-open problem ``quadratic
deformation of the game exterior algebra.''  The implementation has the exterior algebra
$\Lambda_{\Z}M$ of a finitely generated game subgroup $M$, plus a checked
integer-valued Clifford surface that accepts hand-supplied quadratic data only
when the game relations are compatible.  The new observation is that there are
two different descent gates.  A quadratic map must descend through the game
relations; a Clifford quotient over a coefficient ring must also keep the
coefficient ring from being accidentally collapsed by the two-sided relation
ideal.  Over $\Z$ both gates kill torsion squares.  Over torsion coefficient
rings the first gate can reopen, but the second gate rules out several tempting
``universal'' torsion lifts.
\end{abstract}

\section{Live boundary}

\SOURCE{} The implemented object is the following.  Given games
$g_1,\dots,g_n$, the module layer forms the subgroup
\[
  M=\langle g_1,\dots,g_n\rangle_{\Z}
\]
of the normal-play game group, represented by generators and integer relation
rows.  \path{src/games/game_exterior/lambda.rs} builds
$\Lambda_{\Z}M$ by starting from the free Grassmann algebra and quotienting by
the exterior ideal generated by those relation rows.  For example, the relation
$2\star=0$ propagates to $2(\star\wedge\uparrow)=0$.

\SOURCE{} \path{src/games/game_exterior/clifford.rs} is deliberately smaller
than the research problem.  It builds an integer Clifford algebra from
hand-supplied diagonal values $q_i=Q(e_i)$ and off-diagonal polar values
$b_{ij}=B(e_i,e_j)$, then accepts the data only if every game relation is null
and polar-radical.  This is a checked engineering surface, not a source of
game-native quadratic data.

\TESTED{} The focused test suite was run during this pass:
\[
  \texttt{cargo test game\_exterior}.
\]
It passed all twelve filtered tests, including the tests that reject
integer-valued torsion squares and accept the documented torsion vanishings.

\section{Quadratic descent}

Let $R$ be a commutative coefficient ring.  A quadratic datum on an abelian
group $M$ is a map $Q\colon M\to R$ with polar form
\[
  B(x,y)=Q(x+y)-Q(x)-Q(y)
\]
such that $B$ is biadditive and $Q(-x)=Q(x)$.  Equivalently, after choosing
generators, the Clifford relations are
\[
  e_i^2=q_i,\qquad e_i e_j+e_j e_i=b_{ij}.
\]
If a relation row $r=\sum_i c_i e_i$ represents zero in $M$, the quadratic data
descends through the relation only if
\[
  Q(r)=0,\qquad B(r,e_j)=0\quad\text{for every }j.
\]
These are exactly the null and polar-radical checks implemented over $\Z$.

\begin{proposition}[torsion over torsion-free targets]\label{prop:tf}
\PROVED{} Let $A$ be a torsion-free abelian target and let $t\in M$ have finite
order.  For every relation-respecting quadratic datum $Q\colon M\to A$, one
has
\[
  Q(t)=0,\qquad B(t,x)=0\quad\text{for every }x\in M.
\]
\end{proposition}

\begin{proof}
If $nt=0$, then $0=B(nt,x)=nB(t,x)$ for every $x$, so torsion-freeness gives
$B(t,x)=0$.  Also $0=Q(nt)=n^2Q(t)$, hence $Q(t)=0$.
\end{proof}

\begin{corollary}[integer game-Clifford torsion split]\label{cor:intsplit}
\PROVED{} Every integer-valued deformation of a mixed subgroup
$M=T\oplus F$, where $T$ is torsion and $F$ is free, is blind to $T$: the
torsion generators remain exterior/nilpotent and are polar-orthogonal to the
free generators.  All nonzero integer quadratic data factors through the free
quotient.
\end{corollary}

This explains the implemented examples.  In $\langle \star,\uparrow\rangle$,
the relation $2\star=0$ forces $Q(\star)=0$ and $B(\star,\uparrow)=0$ over
$\Z$, while $Q(\uparrow)$ can still be chosen freely.

\section{Clifford-faithful descent is stronger}

For a pure quadratic map, the relation $nt=0$ imposes $n^2Q(t)=0$.  For a
Clifford algebra quotient, that is not the whole story.  The relation vector
$ne_t$ is placed in a two-sided ideal.  Multiplying it by generators can produce
coefficient scalars, so a quotient may accidentally kill part of $R$.

\begin{proposition}[the coefficient-faithfulness obstruction]\label{prop:faith}
\PROVED{} Let $R$ be a commutative coefficient ring and suppose a
Clifford-type quotient imposes the module relation $ne_t=0$ while keeping the
coefficient map $R\to\Cl$ injective.  Then
\[
  nQ(t)=0,\qquad nB(t,x)=0\quad\text{for every }x\in M.
\]
\end{proposition}

\begin{proof}
Since $ne_t$ is zero in the quotient, so is $(ne_t)e_t=nQ(t)$.  If the
coefficient ring injects, this scalar was already zero in $R$.  Similarly
\[
  (ne_t)e_x+e_x(ne_t)=nB(t,x)
\]
is a scalar killed by the relation ideal, so coefficient-faithfulness forces
it to be zero in $R$.
\end{proof}

\begin{remark}
Over $R=\Z$, Proposition~\ref{prop:faith} gives the same visible answer as
Proposition~\ref{prop:tf}: torsion squares vanish.  Over torsion rings it is
strictly sharper.  For a $2$-torsion game $t$, the formal quadratic descent
condition only asks $4Q(t)=0$, but a faithful Clifford quotient asks
$2Q(t)=0$.
\end{remark}

\subsection{The Z/4 Brown lift is not a faithful square quotient}

Take $M=\langle\star\rangle\cong\Z/2$ and try $R=\Z/4$ with $Q(\star)=1$.
As a bare quadratic map, this passes $Q(2\star)=4=0$ in $\Z/4$.  But the
Clifford relation row is $2e_\star=0$, and
\[
  (2e_\star)e_\star=2e_\star^2=2.
\]
Thus the two-sided relation ideal kills the scalar $2$.  The quotient has
silently collapsed the coefficient ring from $\Z/4$ toward characteristic
$2$.  This is not a faithful $\Z/4$-valued deformation.

\SOURCE{} This does not say that $\Z/4$-valued quadratic refinements are
invalid.  They are exactly the Brown-invariant category implemented in
\path{src/forms/char2/brown.rs}: maps $q\colon V\to\Z/4$ on an $\F_2$-space
with
\[
  q(x+y)=q(x)+q(y)+2b(x,y).
\]
That is a genuine quadratic-module target and a real mod-$8$ invariant.  The
point here is narrower: a Brown refinement is not automatically a faithful
Clifford square relation $e_x^2=q(x)$ after also imposing the game-module
relation $2e_x=0$.  The Brown route belongs to the quadratic-module/Form wing;
it does not by itself deform \path{GameExterior}'s algebra without changing the
coefficient ring seen by the quotient.

\subsection{The F2 central-square escape is real but small}

With $R=\F_2$, the same one-generator subgroup can carry the nonzero datum
$Q(\star)=1$ without killing a coefficient scalar: $2=0$ already in $R$.
The resulting algebra is a central-square characteristic-$2$ Clifford algebra,
not a Grassmann algebra.

For a general subgroup $M$, the canonical mod-$2$ reduction gives
$V=M/2M$.  The tautological construction
\[
  R=\Sym_{\F_2}(V),\qquad Q(x)=\bar{x}\in V\subset R,\qquad B=0
\]
is relation-compatible and coefficient-faithful after scalar extension to
characteristic $2$.  It is also too formal to settle the open problem: it has
zero polar form, all generators commute, and the coefficient ring is just the
free polynomial record of the mod-$2$ game class.  It records torsion instead
of explaining it.

\section{Natural positive sources that do not solve torsion}

There is a useful family of positive examples.  If
\[
  \phi\colon M\to A
\]
is an additive game invariant into a commutative torsion-free ring, then
\[
  Q_\phi(x)=\phi(x)^2,\qquad
  B_\phi(x,y)=2\phi(x)\phi(y)
\]
is a natural quadratic datum.  It automatically respects game relations.

\SOURCE{} Two implemented game invariants supply instances of this pattern:

\begin{itemize}[leftmargin=2em]
\item The thermographic mean value is additive in
  \path{src/games/thermography.rs}; this pass ran
  \(\texttt{cargo test mean\_value\_is\_additive}\), which passed.
\item Atomic weight is additive on all-small games in
  \path{src/games/atomic_weight.rs}.  On subgroups where it lands in the
  integer-valued slice, it can be squared as above; this pass ran
  \(\texttt{cargo test atomic\_weight\_is\_additive}\), which passed.
\end{itemize}

For example, on the all-small subgroup generated by $\uparrow$ and $\star$,
atomic weight gives
\[
  \aw(\uparrow)=1,\qquad \aw(\star)=0.
\]
Therefore $Q_{\aw}(\uparrow)=1$ while $Q_{\aw}(\star)=0$ and
$B_{\aw}(\star,\uparrow)=0$.  This is genuinely game-native on the free
all-small direction, and it reproduces the mixed-subgroup split of
Corollary~\ref{cor:intsplit}.  It does not produce nonzero torsion squares.

\section{The nimber core is the one non-tautological torsion source}

The finite nimber subgroup is different because it carries extra structure not
available on arbitrary games: nim-addition, Turning-Corners multiplication, and
trace.  Existing Gold/Arf work uses this to build characteristic-$2$ quadratic
forms
\[
  Q_a(x)=\operatorname{Tr}(x\,x^{2^a})
\]
on finite nimber fields.  This is a real non-tautological torsion quadratic
source, but it lives exactly on the field-like impartial core where the scalar
story already applies.  It does not extend the game exterior algebra over
general partizan games; it confirms the boundary stated in
\path{src/games/partizan.rs}: arbitrary games form an abelian group, not a
ring.

\section{Historical local boundary (superseded globally)}

\RESOLVED{} At the time of this pass, I could not find a non-tautological
torsion square for arbitrary game subgroups.  The alternatives then separated
were:

\begin{enumerate}[leftmargin=2em]
\item Torsion-free coefficient targets.  These are obstructed by
  Proposition~\ref{prop:tf}; torsion games are forced into the exterior
  radical.
\item Torsion coefficient targets with only formal quadratic descent.  These
  can pass $Q(nt)=0$ while failing coefficient-faithful Clifford descent, as
  the $\Z/4$ example shows.
\item Characteristic-$2$ central-square targets.  These can keep nonzero
  torsion squares, but the canonical construction through $M/2M$ is
  tautological and has zero polar form.
\item Nimber-field targets.  These are genuinely game-native through
  Turning-Corners multiplication and trace, but only on the impartial
  field-like core.
\item Additive thermographic or atomic-weight sources.  These are game-native
  and useful on free directions, but every torsion element maps to zero.
\end{enumerate}

At that local stage, the remaining shape looked sharper than ``find a metric.''
A solving construction appeared to need a game-built coefficient target and a
square operation that:

\begin{itemize}[leftmargin=2em]
\item survives the coefficient-faithfulness test of
  Proposition~\ref{prop:faith};
\item is not merely the tautological polynomial ring on $M/2M$;
\item does not factor through an additive invariant into a torsion-free ring;
\item reaches beyond the nimber field core.
\end{itemize}

The later divisibility theorem closes all five commutative-scalar alternatives:
adjoining the required short-game roots kills every torsion square and polar
value before the coefficient characteristic matters.  A game-native
noncommutative or directed structure whose square remembers first-/second-player
asymmetry would instead change the mathematical object and belongs to the
separate outcome-semantics problem.

\section{Implementation consequence}

\PROVED{} The current integer implementation is aligned with the obstruction:
over $\Z$, null and polar-radical relation checks are enough to reject the
torsion squares considered here.  However, a future
generalization of \texttt{GameClifford} to torsion coefficient rings should not
reuse only those checks.  It must also check the scalar intersection of the
two-sided relation ideal, or at least the necessary conditions
\[
  nQ(t)=0,\qquad nB(t,x)=0
\]
for every visible torsion relation $nt=0$.  Otherwise a datum can look
quadratic while the quotient silently changes the coefficient ring.

\end{document}
